\chapter{Radiative and atmospheric forcing}\label{ch:forcing}
\section{Prescribed atmosphere}
The calculation supports a frozen reference atmosphere and a time-dependent prescribed atmosphere. The frozen route uses the midlatitude equinox reference independently of the geographic coordinates used for solar geometry. Changing latitude in a frozen run therefore changes illumination without claiming that the stored reference background follows the new location. The dynamic route evaluates MSIS-00 through the pinned pymsis interface at the requested position and UTC time \citep{picone2002}.

The implementation obtains temperature and ordinary neutral-species densities from MSIS. Total $M$ is the sum of the seven ordinary neutral densities, excluding anomalous oxygen; unavailable trace entries contribute zero to that sum while retaining their availability masks. Densities returned in $\mathrm{m^{-3}}$ are converted to $\mathrm{cm^{-3}}$ by $10^{-6}$. Chemical reservoir fractions are prescribed as $[\oxygen]=0.21M$, $[\mathrm{N_2}]=0.78M$ and $[\mathrm{CO_2}]=405\times10^{-6}M$. These are the model conventions, not a claim that every native MSIS species is used directly in the reaction network.

Water and molecular-hydrogen mixing-ratio profiles follow the historical reference transcription from \citet{brasseur2005}. The external ozone profile also uses that reference prescription because a versioned external CMAM ozone asset was not available. These are approximate prescribed inputs. In particular, external ozone is not a date-dependent observed climatology. Within the chemical domain, ozone and atomic oxygen are supplied by the evolving chemical state; outside it, the external-column conventions provide the attenuation background.

The campaign uses $F_{10.7}=150$, $\overline F_{10.7}=150$ and daily $A_p=4$. The activity interface distinguishes previous-day solar flux and its 81-day average and accepts explicit UTC drivers. It does not retrieve historical activity automatically. Consequently the seasonal cases do not include measured seasonal changes in space weather. Their background differences arise from the prescribed empirical-model response with fixed activity settings.

Background snapshots are computed on a 300-s grid and linearly interpolated in temperature and density. This is a prescribed interpolation scheme, not an integration of atmospheric dynamics. Updating $M(t)$ changes termolecular fluxes, collisional losses and solar attenuation. It does not automatically add a dilution term to an evolving concentration, nor does it move chemical constituents between boxes. The scope of that approximation is discussed in Chapter~\ref{ch:discussion}.

\section{UV/VUV frequencies and evolving opacity}
The ultraviolet/vacuum-ultraviolet calculation retains a historical 125-element spectral backbone from the 2017 asset, including its duplicated tabulated wavelength and Lyman-alpha element. The incident values are photon irradiances per spectral element; the calculation is a discrete sum, not a continuous-spectrum integral requiring an additional wavelength-bin width. This inherited numerical convention is documented in the data provenance rather than attributed to a modern solar-spectrum release.

For absorber $a$ at spectral element $m$, the solar optical depth is assembled from shell densities and path lengths:
\begin{equation}\tau_m(z,t)=\sum_a\sigma_{a,m}(T)\sum_s n_{a,s}(t)\ell_s(z,\chi(t)).
\label{eq:tauUV}\end{equation}
Lengths are converted consistently to centimetres. Illuminated flux is $F_m(z,t)=F_m^0\exp[-\tau_m(z,t)]$; a ray intersecting the solid Earth receives exactly zero. A parent frequency has the discrete form
\begin{equation}J_a(z,t)=\sum_m F_m(z,t)\sigma_{a,m}(T),\label{eq:Jsum}\end{equation}
with the appropriate retained yield included for product-specific quantities.

O and ozone opacity inside 50--100~km changes with the temporal state. The UV frequencies are therefore recomputed using current chemical concentrations rather than imposed as one frozen altitude profile. Oxygen, nitrogen and exterior-column prescriptions supply the other absorbers. The shells extend from the lower column to the fixed 150-km top boundary; radiation above that top is not represented as an additional model atmosphere.

Historical JPL18 water/peroxide cross-section tables and the temperature polynomial supplement the inherited spectral data \citep{jpl18}. A documented wavelength transcription correction uses 189~nm where the historical table printed a sequentially misplaced 199~nm row. Later evaluation material corroborates that correction but does not replace the historical coefficient set. The two retained water channels follow the frozen reduced branching approximation, rather than being described as a complete modern quantum-yield database.

The eleven stored forcing quantities distinguish total oxygen/ozone photolysis, Hartley excitation, SRC, Lyman-alpha, peroxide and two water channels, and the three oxygen-band excitations. Their separate definitions permit the parent/product bookkeeping in Eq.~\eqref{eq:partition}. All are frequencies in $\mathrm{s^{-1}}$.

\section{Oxygen-band excitation}
The A0 baseline uses all 430 A-band, 320 B-band and 835 infrared-atmospheric-band transitions selected from the authorized HITRAN2016 export. Ordinary HITRAN2016 line fields are retained. Temperature-dependent strengths use the frozen TIPS-2017 partition sums; the line shape is classic Voigt \citep{hitran2016,tips2017}. This is a complete pragmatic baseline for the selected bands, not a statement that Voigt is the most detailed available oxygen spectroscopy.

For line $i$, let $S_i(T)$ be the line strength and $V_i(\tilde\nu;T,p)$ its normalized Voigt profile. The molecular absorption cross section is
\begin{equation}\sigma_{\mathrm{O_2}}(\tilde\nu)=\sum_i S_i(T)V_i(\tilde\nu),\qquad
\int_{-\infty}^{\infty}V_i(\tilde\nu)\,d\tilde\nu=1.\label{eq:Voigt}\end{equation}
Doppler and collisional widths enter through the isotope mass, temperature and partial-pressure conventions. Numerical moment/far-wing approximations are used only within the independently audited transfer strategy; they do not remove lines from the band.

The solar source is the fixed Wehrli 1985 extraterrestrial spectrum \citep{wehrli1985}. It is linearly interpolated in wavelength at each transition centre and converted from energy irradiance to photon irradiance per wavenumber using the photon energy and the wavelength--wavenumber Jacobian. Extending that historical spectrum to the three bands is an explicit reconstruction choice. No activity-dependent near-infrared solar variability or unimplemented Earth--Sun distance rescaling is introduced.

The excitation frequency is obtained schematically as
\begin{equation}g_k(z,t)=\sum_{i\in k}\int F_i^0\,S_i(T_z)V_i(\tilde\nu;T_z,p_z)
\exp[-\tau_k(\tilde\nu,z,t)]\,d\tilde\nu.\label{eq:gfactor}\end{equation}
The incoming ray traverses spherical shells; opacity uses the conditions of each shell while the excitation cross section belongs to the target level. This distinction matters when pressure and temperature vary along a long twilight path. The approximate constancy of solar irradiance across an individual line is part of the line-centre solar-source treatment.

For the infrared atmospheric band, historical O$_2$--air collision-induced absorption (CIA) is included in attenuation. The adopted data derive from the 253/273/296-K measurements of \citet{mate1999}. With coefficient $C_{\mathrm{air}}$, the absorption factor uses the historical normalization
\begin{equation}\alpha_{\mathrm{CIA}}(\tilde\nu,T)=C_{\mathrm{air}}(\tilde\nu,T)
[\oxygen]\bigl([\oxygen]+[\mathrm{N_2}]\bigr).\label{eq:CIA}\end{equation}
The reader supplies consistent density units. CIA is not counted as chemical production of $\Delta$. The excitation source remains the monomeric transition sum. A separate O$_2$--O$_2$ term is not added on top of this atmospheric-mixture product. Temperature interpolation and the outside-range envelope retain the historical scope of the data rather than claiming new coefficients.

\section{Advanced A-band sensitivity and baseline choice}
The A1 sensitivity changes exactly 91 principal magnetic-dipole A-band transitions to the validated Drouin speed-dependent Voigt treatment, leaving 339 transitions classic Voigt and omitting line mixing/Galatry. Advanced line-shape physics can matter in precision spectroscopy \citep{drouin2017}; whether it materially changes the outputs here is a separate downstream question.

The temporal comparison re-solves the four noon fast variables for each radiative model while retaining atmosphere, native O/O$_3$/H policy, chemistry and solver controls. Maximum relevant dawn differences are approximately $0.000122\%$ for ozone and $0.233878\%$ for $\Delta$. These are direct output sensitivities of the implemented temporal model, not estimates of absolute spectroscopic accuracy. They support retaining A0 for the final campaign. Exact line mixing and rare-isotopologue Galatry work remain optional for a future precision objective; they are not required to interpret the current figures.

\section{Solar geometry and solid-Earth shadow}
Datetime, latitude and east-positive longitude determine geometric SZA using the documented NOAA fractional-year Fourier equations \citep{noaa}. Datetimes must carry a timezone and are normalized to UTC; geographic latitude lies in $[-90,90]$ degrees and longitude in $[-180,180]$ degrees. The fractional year includes the UTC hour and uses 366 days in leap years. Declination $\delta$ and equation of time $E$ lead to a true solar hour angle $H_s$ and
\begin{equation}\cos\chi=\sin\phi\sin\delta+\cos\phi\cos\delta\cos H_s.\label{eq:SZA}\end{equation}
The implementation is a reproducible Fourier approximation, not a precision ephemeris. It does not apply civil-timezone assumptions, atmospheric refraction or the apparent solar-disk correction used in some sunrise definitions.

Solid-Earth shadow is decided by the atmospheric ray geometry, not by setting every level dark when $\chi>90^\circ$. With Earth radius $R=6370$~km and level radius $r=R+z$, an incoming ray below the local horizon is blocked if its impact parameter $r\sin\chi<R$. The limiting zenith angle is
\begin{equation}\chi_{\mathrm{tan}}(z)=180^\circ-\arcsin\left(\frac{R}{R+z}\right).\label{eq:tangent}\end{equation}
The numerical geometry includes explicitly scoped grazing-roundoff handling. Shadowed solar frequencies are exactly zero, while non-solar chemical sources remain active. Polar day and night follow from the same SZA calculation; a ground-level sunrise is not forced in a case where the geometry has none.

\section{Dynamic forcing tables and output scope}
NIR atmosphere snapshots are taken every 3600~s as a subset of the native 300-s background grid. At the actual SZA, rates interpolate between adjacent atmospheric snapshots. Angular interpolation is independently audited at direct midpoints; no frozen-NIR replacement is silently used for dynamic-background runs. UV uses the current state/background during RHS evaluation.

The forcing is reproducible within the prescriptions, but it is not a historical reconstruction of the Sun and atmosphere on every simulated date. The near-infrared solar spectrum is fixed, activity is quiet unless supplied, and several mixing-ratio profiles remain prescribed. These distinctions are necessary when interpreting the date and latitude experiments as responses of the complete reduced model.
